\documentclass[a4paper,11pt]{article}
\usepackage{exam-style-341}
\examinehead{341 农业综合知识三 · 网络技术分章练习}{第03章 参考答案与解析}
\begin{document}

\unitanswerscover{341 农业综合知识三 · 网络技术分章练习}{第03章参考答案}{物理层：逐题答案与解析；PCM、奈氏准则与香农公式计算题给出完整步骤与单位。}

\answerpart{一、单项选择题（10 题，每题 3 分，共 30 分）}

\q A。\begin{analysis}物理层考虑的是在连接各种计算机的传输媒体上传输数据比特流，其传输单位是比特（教材 2.1，印刷页 42）。帧、分组、报文分别属于数据链路层、网络层和应用层。\end{analysis}

\q A。\begin{analysis}在使用时间域波形表示数字信号时，代表不同离散数值的基本波形称为码元；使用二进制编码时只有两种码元，分别代表 0 和 1（教材 2.2.2，印刷页 44—45）。\end{analysis}

\q B。\begin{analysis}码元传输速率又称波特率，单位是 Baud（码元/秒），表示每秒传输多少个码元（教材 2.2.2，印刷页 45）。bit/s 是比特率单位，Hz 是频率单位。\end{analysis}

\q C。\begin{analysis}每个码元携带的信息量 $=\log_2 V$，$V$ 为码元种类数。16 种码元时 $\log_2 16=4$ bit（教材 2.2.2）。\end{analysis}

\q C。\begin{analysis}奈氏准则：理想低通信道的最高码元传输速率为 $2W$ Baud，$W=4$ kHz 时为 $2\times4000=8000$ Baud（教材 2.2.3，印刷页 46）。\end{analysis}

\q A。\begin{analysis}香农公式为 $C=W\log_2(1+S/N)$，其中 $W$ 为信道带宽，$S/N$ 为信噪比（教材 2.2.3，印刷页 47）。B 是奈氏准则的变形，C 底数错误，D 与公式无关。\end{analysis}

\q A。\begin{analysis}PCM 包括采样、量化、编码三个步骤：先按采样定理对模拟信号采样，再把采样值量化为有限个等级，最后用二进制编码表示（教材 2.5 及通信原理通用内容）。\end{analysis}

\q B。\begin{analysis}根据奈奎斯特采样定理，采样频率应不低于信号最高频率的 2 倍，才能无失真地恢复原信号（教材 2.5；通用教材）。\end{analysis}

\q A。\begin{analysis}FDM（频分复用）按频率划分信道，所有用户在同样的时间占用不同的频带；TDM 按时间划分；STDM 按需分配时隙；CDMA 用正交码区分用户（教材 2.4.1，印刷页 56—59）。\end{analysis}

\q C。\begin{analysis}光纤传输的是光信号而不是电信号，具有传输带宽大、抗电磁干扰、保密性好、损耗小等优点；C 说“传输的是电信号”错误（教材 2.3.1，印刷页 48—52）。\end{analysis}

\answerpart{二、多项选择题（3 题，每题 4 分，共 12 分）}

\q ABC。\begin{analysis}数据通信系统模型由源系统（源点、发送器）、传输系统和目的系统（接收器、终点）组成；路由器是网络层的分组交换设备，不属于该模型的基本组成部分（教材 2.2.1，印刷页 43）。\end{analysis}

\q ABC。\begin{analysis}最基本的带通调制方法有调幅（ASK）、调频（FSK）、调相（PSK），以及正交振幅调制（QAM）；PCM 属于把模拟信号数字化的编码方法，不是带通调制（教材 2.2.2，印刷页 45—46）。\end{analysis}

\q ABC。\begin{analysis}信道复用技术包括频分复用 FDM、时分复用 TDM、统计时分复用 STDM、波分复用 WDM、码分复用 CDMA 等；CSMA/CD 是媒体访问控制协议，不属于信道复用（教材 2.4，印刷页 56—63）。\end{analysis}

\answerpart{三、判断题（6 题，每题 2 分，共 12 分）}

\q ×。改正：只有二进制码元时波特率与比特率数值相等；一般情况下比特率 $=$ 波特率 $\times\log_2 V$（$V$ 为码元种类数）。\begin{analysis}码元传输速率（波特率）与信息传输速率（比特率）概念不同，一个码元可以携带多个比特，因此二者不一定相等（教材 2.2.2，印刷页 45）。\end{analysis}

\q √。\begin{analysis}奈氏准则指出，在带宽为 $W$ 的理想低通信道中，最高码元传输速率为 $2W$ Baud；超过此上限会出现严重码间串扰，使接收端无法正确判决（教材 2.2.3，印刷页 46）。\end{analysis}

\q √。\begin{analysis}香农公式 $C=W\log_2(1+S/N)$ 表明，信道的带宽或信噪比越大，信息的极限传输速率越高（教材 2.2.3，印刷页 47）。\end{analysis}

\q √。\begin{analysis}采样定理要求采样频率不低于信号最高频率的 2 倍，否则会发生频谱混叠，无法无失真恢复原模拟信号（教材 2.5；通用教材）。\end{analysis}

\q ×。改正：光纤传输的是光信号，因此抗电磁干扰能力强，不易受外界电磁干扰。\begin{analysis}光纤利用光的全反射原理传输光信号，具有抗电磁干扰、保密性好等优点（教材 2.3.1）。\end{analysis}

\q ×。改正：统计时分复用按需动态分配时间片，用户有数据时才占用时隙，因此线路利用率比固定时隙的时分复用高。\begin{analysis}固定时隙的 TDM 在用户无数据时隙会浪费；STDM 按需分配，可提高线路利用率（教材 2.4.1）。\end{analysis}

\answerpart{四、计算题（2 题，每题 5 分，共 10 分）}

\q 参考答案：
（1）采样频率 $f_s=2\times f_{\max}=2\times4\ \text{kHz}=8\ \text{kHz}$（1 分）。
（2）量化级数 $N=256$，每个采样值所需编码位数 $n=\log_2 N=\log_2 256=8$ bit（2 分）。
（3）数据率 $=f_s\times n=8000\ \text{次/s}\times8\ \text{bit/次}=64000\ \text{bit/s}=64\ \text{kb/s}$（2 分）。
\begin{analysis}若量化级数改为 128，则 $n=\log_2 128=7$ bit，数据率为 $8000\times7=56$ kb/s，这正是 2022 年真题中“128 个量化级、8 kHz 采样”题的答案思路。计算时注意单位换算：$1\ \text{kb/s}=1000\ \text{bit/s}$（通信领域）。\end{analysis}

\q 参考答案：
（1）信噪比 $30\ \text{dB}$ 对应 $S/N=10^{30/10}=10^{3}=1000$（1 分）；
香农极限 $C=W\log_2(1+S/N)=4000\times\log_2(1+1000)\approx4000\times9.97=39880\ \text{bit/s}\approx39.9\ \text{kb/s}$（2 分）。
（2）采用 16 种码元时，按奈氏准则：
最高码元速率 $=2W=2\times4000=8000$ Baud；
最高信息速率 $=2W\log_2 V=8000\times\log_2 16=8000\times4=32000\ \text{bit/s}=32\ \text{kb/s}$（1.5 分）。
（3）因为 $32\ \text{kb/s}<39.9\ \text{kb/s}$，即低于香农公式给出的极限信息传输速率，所以在该信道上采用 16 种码元可以实现 32 kb/s 的传输（0.5 分）。
\begin{analysis}奈氏准则只考虑带宽和码元种类，给出无噪声条件下的码元速率上限；香农公式考虑带宽和信噪比，给出有噪声信道的信息速率上限。两者从不同角度限制传输速率，实际系统必须同时满足两个约束。\end{analysis}

\answerpart{五、简答题（2 题，每题 8 分，共 16 分）}

\q 参考答案要点：
（1）物理层的主要任务（3 分）：在连接各种计算机的传输媒体上传输数据比特流，尽可能屏蔽掉传输媒体和通信设备的差异，使数据链路层感觉不到这些差异（教材 2.1，印刷页 42）。
（2）接口的四个特性（2 分）：机械特性——接口所用接线器的形状、尺寸、引脚数目和排列等；电气特性——接口电缆各条线上出现的电压范围；功能特性——某条线上某一电平的电压表示何种意义；过程特性——不同功能的各种可能事件的出现顺序。
（3）传输媒体的分类（3 分）：分为导引型传输媒体和非导引型传输媒体两类。导引型包括双绞线、同轴电缆和光纤；非导引型包括无线电波、微波、红外线等（教材 2.3，印刷页 48—55）。

\q 参考答案要点：
（1）奈氏准则（3 分）：在带宽为 $W$（Hz）的理想低通信道中，为避免码间串扰，最高码元传输速率为 $2W$ Baud；若码元有 $V$ 种离散状态，则最高信息传输速率为 $C_{\max}=2W\log_2 V$。它只考虑信道带宽和码元种类，不考虑噪声，适用于理想低通信道（教材 2.2.3，印刷页 46）。
（2）香农公式（3 分）：$C=W\log_2(1+S/N)$，其中 $W$ 为信道带宽，$S/N$ 为信噪比（$S$ 为信号平均功率，$N$ 为噪声平均功率）。它表明信道的极限信息传输速率取决于带宽和信噪比，适用于有随机噪声的信道（教材 2.2.3，印刷页 47）。
（3）区别（2 分）：奈氏准则限制的是码元传输速率，说明提高码元速率要受带宽限制；香农公式限制的是信息传输速率，说明在噪声条件下即使带宽固定，也可通过提高信噪比或增加每码元比特数来提高信息速率，但存在上限。二者分别从“码间串扰”和“噪声”两个角度给出限制，实际系统需同时满足。

\answerpart{六、综合题（1 题，共 20 分）}

\q 参考答案要点：
（1）（5 分）采样频率 $f_s=2\times1\ \text{kHz}=2\ \text{kHz}$；量化级数 64，编码位数 $n=\log_2 64=6$ bit；每路数据率 $=2000\times6=12000\ \text{bit/s}=12\ \text{kb/s}$。
（2）（5 分）FDM 总带宽 $=8\times4\ \text{kHz}+7\times1\ \text{kHz}=32+7=39\ \text{kHz}$（相邻 8 路之间有 7 个保护带）。
（3）（5 分）TDM 每帧携带 $8\times6=48$ bit 数据，加 1 位同步位共 $49$ bit；帧速率为 2000 帧/s，故线路数据率 $=49\times2000=98000\ \text{bit/s}=98\ \text{kb/s}$。
（4）（5 分）信噪比 $20\ \text{dB}$ 对应 $S/N=10^{20/10}=100$；香农极限 $C=W\log_2(1+S/N)=10000\times\log_2 101\approx10000\times6.66=66600\ \text{bit/s}\approx66.6\ \text{kb/s}$。因为 $98\ \text{kb/s}>66.6\ \text{kb/s}$，超过香农极限，所以第（3）问的 TDM 数据率不能在该信道上可靠传输；需增大带宽、提高信噪比、降低采样率或压缩数据。
\begin{analysis}本题把 PCM、FDM、TDM 与香农公式串起来，是 341 网络部分计算题的典型综合形式。易错点：一是 FDM 保护带个数为路数减 1；二是 TDM 帧速率等于采样频率（每帧传一路一个采样值）；三是香农公式中 $S/N$ 是功率比，分贝值要先换算 $S/N=10^{\text{dB}/10}$。\end{analysis}

\end{document}
